\documentclass[11pt,a4paper]{article}

\usepackage[margin=1in]{geometry}
\usepackage{amsmath,amssymb,amsthm,mathtools}
\usepackage[T1]{fontenc}
\usepackage{lmodern}
\usepackage{microtype}
\usepackage{enumitem}
\usepackage{hyperref}
\usepackage{xcolor}

\begin{document}

\section{Kernel Ridge Regression}

In this section we introduce the regression problem that will be used throughout the project. Assume we have a dataset $(x_i,y_i)$ for $i=1,\dots,n$ (this can be e.g. $n$ measurements), where $x_i\in\mathbb{R}^d$ are inputs and $y_i\in\mathbb{R}$ are outputs. We would like to use these data to create a model for the relationship between the $x$ and $y$. In ordinary linear regression one searches over linear functions of the input variables. Kernel ridge regression instead searches over a reproducing kernel Hilbert space $\mathcal{H}$, which allows nonlinear functions to be represented through a kernel function.

Let
\[
k:\mathbb{R}^d \times \mathbb{R}^d \to \mathbb{R}
\]
be a positive definite kernel function. For the training data we define the kernel matrix
\[
K \in \mathbb{R}^{n \times n}, \qquad K_{ij} = k(x_i,x_j).
\]
The entry $K_{ij}$ measures the similarity between the data points $x_i$ and $x_j$ according to the chosen kernel function. The matrix $K$ induced by $k$ is symmetric positive semidefinite.

The kernel trick allows us to formulate this as
\[
\min_{f\in\mathcal H}\|f(X)-y\|_2^2+\lambda\|f\|_{\mathcal H}^2,
\]
which is also
\[
\min_{f\in\mathcal{H}}\sum_{i=1}^n (f(x_i)-y_i)^2+\lambda\|f\|_{\mathcal{H}}^2,
\]
where $\lambda > 0$ is the regularization parameter. The first term measures the error on the training data while the second term penalizes functions with large RKHS(reproducing kernel hilbert space) norm. Here $\lambda$ controls the trade off between fitting the data and overfitting.

By the representer theorem, the minimizer has the form
\[
f(x) = \sum_{j=1}^n \alpha_j k(x,x_j),
\]
for some coefficient vector
\[
\alpha = (\alpha_1,\ldots,\alpha_n) \in \mathbb{R}^n.
\]
For the training points we have
\[
f(x_i) = \sum_{j=1}^n \alpha_j k(x_i,x_j),
\]
so in vector form
\[
f(X) = K\alpha.
\]
Substituting this into the objective gives
\[
\min_{\alpha\in\mathbb R^n}\|K\alpha-y\|_2^2+\lambda\alpha^TK\alpha.
\]
Expand
\[
\|K\alpha-y\|_2^2=(K\alpha-y)^T(K\alpha-y).
\]
Then differentiate
\[
(K\alpha-y)^T(K\alpha-y)+\lambda\alpha^TK\alpha
\]
with respect to $\alpha$. This gives the linear system
\[
(K+\lambda I)\alpha=y.
\]
Therefore the coefficient vector is
\[
\alpha=(K+\lambda I)^{-1}y.
\]
The prediction for a new point $x$ is then
\[
\hat{f}(x)=\sum_{j=1}^n \alpha_j k(x,x_j).
\]

The regularization term is important numerically. Since $K$ can have very small eigenvalues or even zero eigenvalues, adding $\lambda I$ shifts all eigenvalues by $\lambda$ so $K+\lambda I$ is positive definite. This makes the system uniquely solvable and improves the conditioning of the problem.

For large datasets, the main computational difficulty is that $K$ is an $n \times n$ matrix. Forming the full kernel matrix requires $O(n^2)$ kernel evaluations and storing it requires $O(n^2)$ memory. Solving the system directly, for example by Cholesky decomposition costs $O(n^3)$ operations. For small $n$ this is fine but for large datasets this quickly becomes infeasible. Therefore in this project we investigate the Nystr\"om approximation of the kernel matrix: instead of forming all of $K$, we only form $l\ll n$ sampled columns (and the corresponding $l\times l$ submatrix) and make a low rank approximation $\tilde K\approx K$. This is a randomized approach for kernel matrices (Drineas and Mahoney (2005) and Gittens and Mahoney (2016)).

\section{Nystr\"om Approximation}

The Nystr\"om method is a way to approximate a large positive semidefinite matrix by using only a smaller part of it. In this project the positive semidefinite matrix is the kernel matrix $K$, but it is useful to first describe the construction for a general positive semidefinite matrix $A$. The aim is to build a low rank approximation $\tilde{A}$ without using all of $A$ in a full direct factorization.

\subsection{General Nystr\"om form}

Following the general formulation in randomized numerical linear algebra, let
\[
X \in \mathbb{R}^{n \times l}
\]
be a test matrix, the Nystr\"om approximation of $A$ with respect to $X$ is
\[
A\langle X\rangle=(AX)(X^TAX)^\dagger(AX)^T.
\]
Here $(X^TAX)^\dagger$ denotes the Moore--Penrose pseudoinverse. The matrix $AX$ contains the information we collect from $A$.

This formula shows why the approximation has rank at most $l$ since the formula is a product involving the $n \times l$ matrix $AX$, so
\[
\operatorname{rank}(A\langle X\rangle) \leq l.
\]
Therefore, when $l \ll n$, the approximation is much cheaper to store and compute with than the full matrix $A$.

There is also a useful geometric representation of this formula. Since $A$ is positive semidefinite we can use $A^{1/2} $. Then
\[
A\langle X\rangle=A^{1/2} P_{A^{1/2}X} A^{1/2},
\]
where $P_{A^{1/2}X}$ is the orthogonal projector onto the column space of $A^{1/2}X$. This has error as
\[
A - A\langle X\rangle=A^{1/2}(I-P_{A^{1/2}X})A^{1/2}
\]
which is precisely the Schur complement of $A$ with respect to $X$.

\subsection{Randomized Nystr\"om approximation}

The randomized Nystr\"om approximation draws a random matrix
\[
\Omega \in \mathbb{R}^{n \times l}
\]
and form
\[
Y = A\Omega.
\]
Then the Nystr\"om approximation is
\[
\hat{A}=Y(\Omega^TY)^\dagger Y^T.
\]

\end{document}